\documentclass[11pt,a4paper]{article}

\usepackage{amsmath, amssymb, amsthm}
\usepackage{geometry}
\usepackage{microtype}
\usepackage{parskip}
\usepackage{booktabs}
\usepackage[colorlinks=true, linkcolor=blue, citecolor=blue, urlcolor=blue]{hyperref}

\geometry{margin=2.2cm}

\title{\textbf{Nonlinear HJB Equations in Finance:\\Selected Examples}}
\author{}
\date{}

\begin{document}

\maketitle
\thispagestyle{empty}
\vspace{-1.5em}

\noindent
In classical derivative pricing, the governing PDEs---Black-Scholes, Heston, CIR---are \emph{linear} in the
unknown $V$, even when their coefficients depend on state variables such as the spot price or stochastic
variance. Linearity is preserved because the replicating hedge is determined exogenously by no-arbitrage
alone. The situation changes fundamentally when an agent optimises over a control. The Hamilton-Jacobi-Bellman
(HJB) equation, which characterises the value function of a stochastic optimal control problem via dynamic
programming, is generically \emph{nonlinear} in the value function: after substituting the optimal control
back into the right-hand side, the maximised Hamiltonian produces terms that are rational or quadratic in
the value function and its derivatives. This endogenous nonlinearity is structurally analogous to the
advection term $(\mathbf{u}\cdot\nabla)\mathbf{u}$ in Navier-Stokes, where the solution itself acts as its
own coefficient. Superposition fails, classical solution theory no longer applies, and the appropriate
framework becomes that of \emph{viscosity solutions}. The examples below illustrate the range of nonlinear
HJB equations that arise in quantitative finance.

\bigskip
\hrule
\bigskip

\subsection*{1. Merton Portfolio Optimisation}

An investor allocates fraction $\pi$ of wealth $W$ to a risky asset with drift $\mu$ and volatility $\sigma$,
maximising expected utility. The HJB equation for the value function $V(W,t)$ is
\[
  \partial_t V + \max_{\pi}\left\{
    \tfrac{1}{2}\sigma^2\pi^2 W^2\, V_{WW} + (\mu\pi W + rW)\,V_W
  \right\} = 0.
\]
Substituting the optimal weight $\pi^* = -(\mu - r)\,V_W / (\sigma^2 W V_{WW})$ yields the reduced PDE
\[
  \partial_t V + rW\,V_W - \frac{(\mu-r)^2}{2\sigma^2}\frac{V_W^2}{V_{WW}} = 0,
\]
where the term $(V_W)^2 / V_{WW}$ is a \emph{rational} nonlinearity in the value function---absent from
any linear pricing PDE.

\subsection*{2. Uncertain Volatility / Robust Control (Black-Scholes-Barenblatt)}

When volatility is unspecified but bounded, $\sigma \in [\sigma_{\min}, \sigma_{\max}]$, the worst-case
pricing problem is a zero-sum stochastic game. The value function satisfies the
Hamilton-Jacobi-Bellman-Isaacs (HJBI) equation
\[
  \partial_t V + \frac{1}{2}\hat{\sigma}\!\left(\tfrac{\partial^2 V}{\partial S^2}\right)^{\!2}
  S^2\,\frac{\partial^2 V}{\partial S^2} + rS\,\frac{\partial V}{\partial S} - rV = 0,
\]
where $\hat{\sigma}(\Gamma) = \sigma_{\max}$ if $\Gamma > 0$ and $\sigma_{\min}$ if $\Gamma < 0$. The
diffusion coefficient thus depends on the \emph{sign of} $\partial^2 V/\partial S^2$, making the effective
volatility endogenous to the solution itself.

\subsection*{3. Optimal Execution (Almgren-Chriss type)}

A trader liquidates inventory $q$ at rate $u$ over $[0,T]$, minimising expected cost plus quadratic
inventory risk. The HJB equation is
\[
  \partial_t V + \min_{u}\!\left\{
    -u\,V_q + \eta\, u^2 + \tfrac{1}{2}\gamma\sigma^2 q^2
  \right\} = 0.
\]
Inserting $u^* = V_q/(2\eta)$ gives
\[
  \partial_t V - \frac{(V_q)^2}{4\eta} + \frac{1}{2}\gamma\sigma^2 q^2 = 0.
\]
The term $(V_q)^2/4\eta$ is \emph{quadratic} in the first derivative of $V$---the precise analogue of
$(\mathbf{u}\cdot\nabla)\mathbf{u}$ in Navier-Stokes.

\subsection*{4. Differential Borrowing and Lending Rates}

When the borrowing rate $r_b$ exceeds the lending rate $r_l$, the relevant rate depends on the sign of the
cash account $X = V - SV_S$. The option value satisfies
\[
  \partial_t V + \tfrac{1}{2}\sigma^2 S^2 V_{SS} + rS\,V_S
  - r_l\max(V - SV_S,\,0) - r_b\min(V - SV_S,\,0) = 0.
\]
Nonlinearity enters through $\max(\cdot,0)$ and $\min(\cdot,0)$, which depend on $V$ and $V_S$.

\subsection*{5. Multi-Asset Market Making (Bergault, Evangelista, Gu\'eant, Vieira)}

Consider a market maker quoting bid and ask prices across $d$ correlated assets with inventory vector
$q\in\mathbb{R}^d$ and reference prices following correlated Brownian motions with covariance matrix
$\Sigma$. Under Model~B (mean-variance objective), the HJB equation reduces via an ansatz to the
Hamilton-Jacobi equation for the value adjustment $\theta(t,q)$:
\[
  0 = \partial_t\theta(t,q) - \frac{\gamma}{2}\,q^\top\Sigma\, q
  + \sum_{i=1}^{d} z^i\left[
    H^{i,b}_0\!\!\left(\frac{\theta(t,q)-\theta(t,q+z^i e_i)}{z^i}\right)
    + H^{i,a}_0\!\!\left(\frac{\theta(t,q)-\theta(t,q-z^i e_i)}{z^i}\right)
  \right],
\]
with terminal condition $\theta(T,q)=0$, where $H^{i,b}_0$ and $H^{i,a}_0$ are the \emph{Hamiltonian
functions} defined by
\[
  H^{i,b}_0(p) = \sup_{\delta}\,\Lambda^{i,b}(\delta)\,(\delta - p),
\]
and analogously for the ask side. The Hamiltonians encode the liquidity-adjusted trade-off between
spread revenue and fill probability. Since $H^{i,b}_0$ and $H^{i,a}_0$ are nonlinear functions evaluated
at \emph{finite differences of $\theta$}, the equation is nonlinear in the value function. The optimal
bid and ask quotes are then recovered as
\[
  \delta^{i,b\,*}_t = \tilde\delta^{i,b\,*}_0\!\left(
    \frac{\theta(t,q_{t-}) - \theta(t,q_{t-}+z^i e_i)}{z^i}
  \right), \qquad
  \delta^{i,a\,*}_t = \tilde\delta^{i,a\,*}_0\!\left(
    \frac{\theta(t,q_{t-}) - \theta(t,q_{t-}-z^i e_i)}{z^i}
  \right).
\]
Because the true $\theta$ is computationally intractable for large $d$, the paper approximates the
Hamiltonians by quadratic functions, reducing the HJB to a matrix Riccati ODE solvable in closed form.
In the ergodic limit, the optimal half-spread and skew for asset $i$ simplify to
\[
  \frac{\breve\delta^{i,a}_t + \breve\delta^{i,b}_t}{2}
  = \frac{1}{2}\sqrt{\gamma}\,z^i\,e_i^\top \Gamma\, e_i + \frac{1}{k^i},
  \qquad
  \frac{\breve\delta^{i,a}_t - \breve\delta^{i,b}_t}{2}
  = -\sqrt{\gamma}\,q_{t-}^\top \Gamma\, e_i,
\]
where $\Gamma = D_+^{-1/2}(D_+^{1/2}\Sigma D_+^{1/2})^{1/2}D_+^{-1/2}$ encodes the interplay between
price risk and liquidity risk across all assets.

\bigskip
\hrule
\bigskip

\noindent\small
\textbf{Summary.} In all five cases, nonlinearity in the value function arises from substituting an
optimal control back into the Hamiltonian. The locus and structure of the nonlinearity varies---rational
($V_W^2/V_{WW}$ in Merton), quadratic in gradients (optimal execution), switching (uncertain volatility,
differential rates), or evaluated at finite differences of $\theta$ (market making)---but the common
feature is that the solution itself determines the operator acting on it, precluding superposition and
requiring viscosity solution theory for rigorous treatment.

\end{document}
